\documentclass[10pt,a4paper]{amsart}
\usepackage[margin=19mm]{geometry}
\usepackage{amsmath,amssymb,mathtools}
\usepackage[scr=rsfs,cal=euler]{mathalfa}
\usepackage{xfrac}
\usepackage[unicode,hidelinks,bookmarks=false]{hyperref}
\newcommand{\ie}{i.e.\ }
\newcommand{\cf}{cf.\ }
\newcommand{\resp}{respectively\ }
\newcommand{\Subsection}{\subsection}
\providecommand{\nobreakdash}{\nobreak\hbox{-}}
\setlength{\emergencystretch}{2em}
\multlinegap0pt
\usepackage[shortlabels]{enumitem}
\usepackage{amssymb}
\usepackage{xcolor}
\usepackage{mathtools}
\usepackage{tikz}
\usepackage{dsfont}
\newtheorem{theorem}{Theorem}
\title{Stable 2--systole bounds in positive scalar curvature}
\author{Douglas Stryker}
\date{Source: arXiv:2604.22106v1 (CC BY 4.0)}
\begin{document}
\maketitle

\noindent\emph{Source and licence.} Stable 2--systole bounds in positive scalar curvature. Version arXiv:2604.22106v1 (CC BY 4.0), distributed by arXiv under the Creative Commons Attribution 4.0 International licence, http://creativecommons.org/licenses/by/4.0/, as recorded in the arXiv licence metadata for this exact version and shown on the abstract page https://arxiv.org/abs/2604.22106v1. Full licence notice: \texttt{licenses/arxiv-2604.22106-CC-BY-4.0.txt}.\par\smallskip

\noindent\textit{Editorial scope.} Counted unit: the article's theorem thm:txs2 (source0.tex lines 150--152). The article's complete original proof of it is delivered with it (source0.tex lines 578--607, one \texttt{proof} environment of 30 lines, which the article places away from the statement). No mathematics is written, repaired or re-derived: the statement and the proof are the article's own lines, in the article's own notation, and no retained line is substituted or re-flowed.  Retained uncounted as the source-local context the proof needs: lines 116--119; lines 156--158.  Nothing was omitted from any retained span, so the package carries no omission record.  No retained line contains a citation, so the wrapper carries no bibliography and no bibliography entry.  No retained line contains a figure environment, an image-inclusion command or drawing code, so no image is delivered and the package declares no asset dependency.  Editorial formatting only, in addition: an AMS \texttt{amsart} preamble that restates the article's own declarations and macros used by the retained lines, namely the environments \texttt{theorem} on separate counters, together with the standard packages those lines need.  The retained environments print in this document as Theorem 1 in order of appearance, whereas the article prints the same items with the numbering of its own full document.  The complete native source of the article as distributed in the arXiv e-print for this exact version is bundled unchanged as \texttt{sources/199072/source0.tex}.
\begin{theorem}\label{thm:gen}
Let $M$ be a $2m$--dimensional closed 2--essential manifold. Let $N$ be a $2k$--dimensional closed enlargeable manifold, or a point. Suppose $M \times N$ is spin, and let $b = b_2(M \times N)$ be the second Betti number. Let $g$ be a Riemannian metric on $M \times N$ with $\mathrm{scal}(g) \geq 1$. Then
\[ \mathrm{stsys}_2(g) \leq C_{b, k, m}. \]
\end{theorem}

\begin{theorem}\label{thm:txs2n}
If $((S^2)^m \times T^n, g)$ satisfies $\mathrm{scal}(g) \geq 1$, then $\mathrm{stsys}_2^{\mathrm{sph}}(g) \leq C_{m,n}$.
\end{theorem}

\begin{theorem}\label{thm:txs2}
If $(S^2 \times T^n, g)$ satisfies $\mathrm{scal}(g) \geq 1$, then $\mathrm{stsys}_2^{\mathrm{sph}}(g) \leq C_n$.
\end{theorem}

\begin{proof}[Proof of Theorem \ref{thm:txs2} and \ref{thm:txs2n}]
We note that Theorem \ref{thm:txs2} follows from the more general Theorem \ref{thm:txs2n}, which we prove now.

If $n$ is odd, then we can take a Riemannian product with a large circle and the conclusion holds. Hence, we assume $n = 2k$.

Let $\pi_l: ((S^2)^m \times T^{2k}, \tilde{g}) \to ((S^2)^m \times T^{2k}, g)$ be the cover corresponding to the $l$--fold cover of each $S^1$ factor in $T^{2k}$. Let $(\pi_l)_* : \Lambda_2((S^2)^m \times T^{2k}) \to \Lambda_2((S^2)^m \times T^{2k})$ be the restriction of the induced homomorphism on homology.

For $j \in \{1, \hdots, m\}$, let $\sigma_j$ be the homology class of points times the 2--sphere in the $j$--th factor:
\[ \{\mathrm{pt}\} \times \cdots \times S^2 \times \cdots \times \{\mathrm{pt}\} \times \{\mathrm{pt}\}. \]
For $l = \binom{2k}{2}$, let $\tau_1, \hdots, \tau_l$ be the homology classes of points times two circles in the $T^{2k}$ factor. Then $\{\sigma_1, \hdots, \sigma_m, \tau_1, \hdots, \tau_l\}$ is a generating set for $\Lambda_2((S^2)^m \times T^{2k})$. Since the covering space and the base have the same topology, we can use the same basis for both. To avoid ambiguity, we will put a subscript on the stable norm to specify which metric we are using.

We note that
\begin{equation}\label{eqn:cover} (\pi_l)_*(\sigma_j) = \sigma_j,\ \ (\pi_l)_*(\tau_j) = l^2\tau_j. \end{equation}
Since $\pi$ is a local isometry, for any $\sigma \in \Lambda_2((S^2)^m \times T^{2k})$, we have
\begin{equation}\label{eqn:norm-increasing} \|\sigma\|_{\tilde{g}} \geq \|(\pi_l)_*(\sigma)\|_g. \end{equation}
Moreover, since $S^2$ is simply connected, there is a constant $C > 0$ so that
\[ \|\sigma_j\|_{\tilde{g}} \leq C \]
for all $j$ and all $l$. Hence, for $l$ sufficiently large, we have
\[ \mathrm{stsys}_2(\tilde{g}) = \inf\{\|\sigma\|_{\tilde{g}} : \sigma \in \mathrm{Span}(\sigma_1, \hdots, \sigma_j) \setminus \{0\}\}. \]
Applying Theorem \ref{thm:gen} to $\tilde{g}$ (with, for example, $M = (S^2)^m \times T^{2k}$ and $N = \{\mathrm{pt}\}$) and using \eqref{eqn:cover} and \eqref{eqn:norm-increasing}, we have
\begin{align*}
\mathrm{stsys}_2^{\mathrm{sph}}(g)
& = \inf\{\|\sigma\|_g : \sigma \in \mathrm{Span}(\sigma_1, \hdots, \sigma_j) \setminus \{0\}\}\\
& = \inf\{\|(\pi_l)_*(\sigma)\|_g : \sigma \in \mathrm{Span}(\sigma_1, \hdots, \sigma_j) \setminus \{0\}\}\\
& \leq \inf\{\|\sigma\|_{\tilde{g}} : \sigma \in \mathrm{Span}(\sigma_1, \hdots, \sigma_j) \setminus \{0\}\}\\
& = \mathrm{stsys}_2(\tilde{g})\\
& \leq C_{m, n},
\end{align*}
as desired.
\end{proof}

\end{document}
